\documentclass[12pt]{article}
\usepackage[utf8]{inputenc}
\usepackage{geometry}
\usepackage{graphicx}
\usepackage{amsmath}
\usepackage{amssymb}

\usepackage{hyperref}

\geometry{margin=1in}

\title{A Sample Research Paper}
\author{John Doe}
\date{\today}

\begin{document}

\maketitle

\begin{abstract}
This is a sample abstract for a research paper. It demonstrates the use of the \texttt{{changes}} package for tracking revisions. The new findings suggest that this approach is effective for collaborative writing.
\end{abstract}

\section{Introduction}
\label{sec:intro}

The quick brown fox jumps over the lazy dog. This is a  paragraph that demonstrates text processing.

\subsection{Background}
\label{subsec:background}

Previous work has been done in this area~\cite{ref2021}. The important contributions include:
\begin{itemize}\itemFirst contribution
 \itemSecond contribution
 \itemThird contribution
\end{itemize}

\section{Methodology}
\label{sec:method}

We used the following equation:
\begin{equation}
 E = mc^2
 \label{eq:energy}
\end{equation}

The old method produces better results.

\subsection{Data Collection}
Data was collected from multiple sources. The \cancel{invalid} entries were removed.

\section{Results}
\label{sec:results}

Our results are shown in Figure~\ref{fig:example}. The data indicates:
\begin{enumerate}\itemFirst result
 \itemSecond result
 \itemThird result
\end{enumerate}

\begin{figure}
 \centering
 \includegraphics[width=0.5\textwidth]{example.png}
 
 
\caption{Example figure}\label{fig:example}\end{figure}

\section{Discussion}
\label{sec:discussion}

We discuss our findings here. The  is no longer valid. The new interpretation provides better insight.

\section{Conclusion}
\label{sec:conclusion}

In conclusion, this paper demonstrates the effectiveness of the \texttt{{changes}} package. Future work will explore additional applications.

\bibliographystyle{plain}
\bibliography{references}

\end{document}